% GENERATED FILE. Edit canonical lessons, then run npm run generate:books.
% canonical-source-hash: fnv1a64:f73d73f4d27c6e99
% canonical-lessons: TA-C327-munnor, TA-C327-vamcam, TA-C327-kutiyeru, TA-C327-netunkalam, TA-C327-purvikam

\chapter{Ancestors}
\label{ch:ta-ancestors}
\hlchaptermodality{327}

\begin{chapteropening}
\textbf{By the end of this chapter:} \emph{I can talk about my ancestors and where they settled.}

Complete the last lesson of chapter 327: I can talk about my ancestors and where they settled.
\end{chapteropening}

\section[\texorpdfstring{\ta{முன்னோர்} (muṉṉōr)}{முன்னோர் (muṉṉōr)}]{\ta{முன்னோர்} (muṉṉōr) — ancestors; forebears}

\label{lesson:TA-C327-munnor}



\begin{quote}
Before the new one: say the Tamil for a period, then the Tamil for a decade.
\end{quote}

\subsection*{You'll want to know: \ta{முன்னோர்}}
\textbf{\ta{முன்னோர்}} — \emph{muṉṉōr} — \textquotedblleft{}ancestors; forebears\textquotedblright{}.

\begin{quote}
\ta{எங்கள்} \ta{முன்னோர்} \ta{இந்த} \ta{ஊரில்} \ta{வாழ்ந்தார்கள்}.
\end{quote}

\begin{quote}
\emph{eṅkaḷ muṉṉōr inta ūril vāḻntārkaḷ.}
\end{quote}

\begin{quote}
Our ancestors lived in this village.
\end{quote}

Plural in sense; \ta{முன்னோர்கள்} is also heard.

\begin{cousinweb}
\textbf{\ta{முன்னோர்}} (\emph{muṉṉōr}) is native Dravidian: those who were before, from \ta{முன்}, \textquotedblleft{}in front, before\textquotedblright{} (DEDR 5020).
\end{cousinweb}

\begin{grammarlens}[title={What you've built}]
One more word for telling what happened.
\end{grammarlens}

\subsection*{Your turn}
\begin{itemize}
\raggedright
  \item \emph{Say it:} \emph{muṉṉōr}
  \item \emph{Say it:} \emph{muṉṉōr}, once more
  \item \emph{Say it:} \emph{eṅkaḷ muṉṉōr inta ūril vāḻntārkaḷ.} — Our ancestors lived in this village.
  \item \emph{Recall:} say the Tamil for a period, then the Tamil for a decade, then say \emph{muṉṉōr} again
\end{itemize}

\begin{tcolorbox}[breakable,colback=teal!4,colframe=teal!35!black,title={Before you move on}]
What does \ta{முன்னோர்} mean? (\textquotedblleft{}Ancestors; forebears\textquotedblright{}.) Say it once more.
\end{tcolorbox}
\section[\texorpdfstring{\ta{வம்சம்} (vamcam)}{வம்சம் (vamcam)}]{\ta{வம்சம்} (vamcam) — a dynasty; a family line}

\label{lesson:TA-C327-vamcam}



\begin{quote}
Before the new one: say the Tamil for a decade, then the Tamil for ancestors.
\end{quote}

\subsection*{You'll want to know: \ta{வம்சம்}}
\textbf{\ta{வம்சம்}} — \emph{vamcam} — \textquotedblleft{}a dynasty; a family line\textquotedblright{}.

\begin{quote}
\ta{அவர்} \ta{ஒரு} \ta{பழைய} \ta{வம்சத்தைச்} \ta{சேர்ந்தவர்}.
\end{quote}

\begin{quote}
\emph{avar oru paḻaiya vamcattaic cērntavar.}
\end{quote}

\begin{quote}
He belongs to an old family line.
\end{quote}

Also a royal line.

\begin{cousinweb}
\textbf{\ta{வம்சம்}} (\emph{vamcam}) is borrowed from Sanskrit \emph{vaṃśa}, \textquotedblleft{}bamboo\textquotedblright{}, and so \textquotedblleft{}family, lineage\textquotedblright{} (CDIAL 11175), a line that grows joint by joint.
\end{cousinweb}

\begin{grammarlens}[title={What you've built}]
One more word for telling what happened.
\end{grammarlens}

\subsection*{Your turn}
\begin{itemize}
\raggedright
  \item \emph{Say it:} \emph{vamcam}
  \item \emph{Say it:} \emph{vamcam}, once more
  \item \emph{Say it:} \emph{avar oru paḻaiya vamcattaic cērntavar.} — He belongs to an old family line.
  \item \emph{Recall:} say the Tamil for a decade, then the Tamil for ancestors, then say \emph{vamcam} again
\end{itemize}

\begin{tcolorbox}[breakable,colback=teal!4,colframe=teal!35!black,title={Before you move on}]
What does \ta{வம்சம்} mean? (\textquotedblleft{}A dynasty; a family line\textquotedblright{}.) Say it once more.
\end{tcolorbox}
\section[\texorpdfstring{\ta{குடியேறு} (kuṭiyēṟu)}{குடியேறு (kuṭiyēṟu)}]{\ta{குடியேறு} (kuṭiyēṟu) — to settle; to move in}

\label{lesson:TA-C327-kutiyeru}



\begin{quote}
Before the new one: say the Tamil for ancestors, then the Tamil for a dynasty.
\end{quote}

\subsection*{You'll want to know: \ta{குடியேறு}}
\textbf{\ta{குடியேறு}} — \emph{kuṭiyēṟu} — \textquotedblleft{}to settle; to move in\textquotedblright{}.

\begin{quote}
\ta{என்} \ta{தாத்தா} \ta{இளம்} \ta{வயதில்} \ta{இந்த} \ta{நகரத்தில்} \ta{குடியேறினார்}.
\end{quote}

\begin{quote}
\emph{eṉ tāttā iḷam vayatil inta nakarattil kuṭiyēṟiṉār.}
\end{quote}

\begin{quote}
My grandfather settled in this city when he was young.
\end{quote}

Past (I …): \ta{குடியேறினேன்}, \emph{kuṭiyēṟiṉēṉ}.

The place takes -\ta{இல்}; \ta{குடி}, home + \ta{ஏறு}, to go up into.

\begin{cousinweb}
\textbf{\ta{குடியேறு}} (\emph{kuṭiyēṟu}) is a Tamil compound built from \ta{குடி}, \textquotedblleft{}house, home, family\textquotedblright{} (DEDR 1655), and \ta{ஏறு}, \textquotedblleft{}to rise, go up into\textquotedblright{} (DEDR 916).
\end{cousinweb}

\begin{grammarlens}[title={What you've built}]
One more word for telling what happened.
\end{grammarlens}

\subsection*{Your turn}
\begin{itemize}
\raggedright
  \item \emph{Say it:} \emph{kuṭiyēṟu}
  \item \emph{Say it:} \emph{kuṭiyēṟu}, once more
  \item \emph{Say it:} \emph{eṉ tāttā iḷam vayatil inta nakarattil kuṭiyēṟiṉār.} — My grandfather settled in this city when he was young.
  \item \emph{Recall:} say the Tamil for ancestors, then the Tamil for a dynasty, then say \emph{kuṭiyēṟu} again
\end{itemize}

\begin{tcolorbox}[breakable,colback=teal!4,colframe=teal!35!black,title={Before you move on}]
What does \ta{குடியேறு} mean? (\textquotedblleft{}To settle; to move in\textquotedblright{}.) Say it once more.
\end{tcolorbox}
\section[\texorpdfstring{\ta{நெடுங்காலம்} (neṭuṅkālam)}{நெடுங்காலம் (neṭuṅkālam)}]{\ta{நெடுங்காலம்} (neṭuṅkālam) — a long time; ages}

\label{lesson:TA-C327-netunkalam}



\begin{quote}
Before the new one: say the Tamil for a dynasty, then the Tamil for to settle.
\end{quote}

\subsection*{You'll want to know: \ta{நெடுங்காலம்}}
\textbf{\ta{நெடுங்காலம்}} — \emph{neṭuṅkālam} — \textquotedblleft{}a long time; ages\textquotedblright{}.

\begin{quote}
\ta{இந்தக்} \ta{கோயில்} \ta{நெடுங்காலமாக} \ta{இங்கே} \ta{இருக்கிறது}.
\end{quote}

\begin{quote}
\emph{intak kōyil neṭuṅkālamāka iṅkē irukkiṟatu.}
\end{quote}

\begin{quote}
This temple has been here for a long time.
\end{quote}

With -\ta{ஆக}, for a long time: \ta{நெடுங்காலமாக}.

\begin{cousinweb}
\textbf{\ta{நெடுங்காலம்}} (\emph{neṭuṅkālam}) joins native \ta{நெடு}, \textquotedblleft{}long\textquotedblright{} (DEDR 3738), with \ta{காலம்}, \textquotedblleft{}time\textquotedblright{}, borrowed from Sanskrit \emph{kāla} (CDIAL 3084).
\end{cousinweb}

\begin{grammarlens}[title={What you've built}]
One more word for telling what happened.
\end{grammarlens}

\subsection*{Your turn}
\begin{itemize}
\raggedright
  \item \emph{Say it:} \emph{neṭuṅkālam}
  \item \emph{Say it:} \emph{neṭuṅkālam}, once more
  \item \emph{Say it:} \emph{intak kōyil neṭuṅkālamāka iṅkē irukkiṟatu.} — This temple has been here for a long time.
  \item \emph{Recall:} say the Tamil for a dynasty, then the Tamil for to settle, then say \emph{neṭuṅkālam} again
\end{itemize}

\begin{tcolorbox}[breakable,colback=teal!4,colframe=teal!35!black,title={Before you move on}]
What does \ta{நெடுங்காலம்} mean? (\textquotedblleft{}A long time; ages\textquotedblright{}.) Say it once more.
\end{tcolorbox}
\section[\texorpdfstring{\ta{பூர்வீகம்} (pūrvīkam)}{பூர்வீகம் (pūrvīkam)}]{\ta{பூர்வீகம்} (pūrvīkam) — ancestral roots; native place}

\label{lesson:TA-C327-purvikam}



\begin{quote}
Before the new one: say the Tamil for to settle, then the Tamil for a long time.
\end{quote}

\subsection*{You'll want to know: \ta{பூர்வீகம்}}
\textbf{\ta{பூர்வீகம்}} — \emph{pūrvīkam} — \textquotedblleft{}ancestral roots; native place\textquotedblright{}.

\begin{quote}
\ta{எங்கள்} \ta{குடும்பத்தின்} \ta{பூர்வீகம்} \ta{ஒரு} \ta{சிறிய} \ta{கிராமம்}.
\end{quote}

\begin{quote}
\emph{eṅkaḷ kuṭumpattiṉ pūrvīkam oru ciṟiya kirāmam.}
\end{quote}

\begin{quote}
Our family comes from a small village.
\end{quote}

The place a family comes from; \ta{பூர்வீக} \ta{வீடு}, the ancestral house.

\begin{cousinweb}
\textbf{\ta{பூர்வீகம்}} (\emph{pūrvīkam}) is usually traced to Sanskrit \emph{pūrvika}, \textquotedblleft{}former, ancient\textquotedblright{} (Monier-Williams), with the Tamil -\ta{அம்} of nouns.
\end{cousinweb}

\begin{grammarlens}[title={What you've built}]
One more word for telling what happened.
\end{grammarlens}

\subsection*{Your turn}
\begin{itemize}
\raggedright
  \item \emph{Say it:} \emph{pūrvīkam}
  \item \emph{Say it:} \emph{pūrvīkam}, once more
  \item \emph{Say it:} \emph{eṅkaḷ kuṭumpattiṉ pūrvīkam oru ciṟiya kirāmam.} — Our family comes from a small village.
  \item \emph{Recall:} say the Tamil for to settle, then the Tamil for a long time, then say \emph{pūrvīkam} again
\end{itemize}

\begin{tcolorbox}[breakable,colback=teal!4,colframe=teal!35!black,title={Before you move on}]
What does \ta{பூர்வீகம்} mean? (\textquotedblleft{}Ancestral roots; native place\textquotedblright{}.) Say it once more.
\end{tcolorbox}
